\documentclass[aps,prd,twocolumn,nofootinbib,superscriptaddress,10pt]{revtex4-2}

\usepackage[utf8]{inputenc}
\usepackage{amsmath, amssymb, amsfonts, amsthm}
\usepackage{graphicx}
\usepackage{hyperref}
\usepackage{physics}
\usepackage{tensor}
\usepackage{xcolor}

\newtheorem{remark}{Remark}

\hypersetup{
    colorlinks=true,
    linkcolor=blue,
    citecolor=red,
    urlcolor=cyan
}

\begin{document}

\title{Topological Origins of the Dark Sector via K3 Compactification:\\ The String Axiverse, UV Consistency, and Vacuum Energy}

\author{Xavier Callens}
\affiliation{SocrateAI Open Lab, Independent Research Node, France}

\date{\today}

\begin{abstract}
We propose a phenomenological ultraviolet (UV) embedding for Fuzzy Dark Matter (FDM) and Dark Energy by mapping the $S_{1,2}$ Picard-Fuchs effective field theory into a Type IIB string compactification on a K3-fibered Calabi-Yau threefold. By identifying the fundamental constituents of the dark sector with the exact topological invariants of a K3 surface, we construct a highly constrained framework. In alignment with recent extra-dimensional resonance models (Lee \& Tsai, Phys. Rev. D, 2026), the integration of the Ramond-Ramond 4-form over the $b_2=22$ homology cycles naturally generates the String Axiverse. Here, the exact-rational $S_{1,2}$ sequence governs the Euclidean D3-brane instanton action, isolating a $\sim 10^{-21}$ eV Kaluza-Klein resonance suitable for FDM. Furthermore, the K3 Euler characteristic ($\chi=24$) provides the necessary geometric condition for integer D3-brane tadpole cancellation, explicitly evading the Swampland, while simultaneously driving the $\alpha'$ corrections of the Large Volume Scenario (LVS) to generate a metastable de Sitter vacuum. Finally, the severe chiral asymmetry dictated by the Hirzebruch signature ($\tau=-16$) yields a stress-energy Trace Anomaly via the Atiyah-Singer Index Theorem, providing a rigorous quantum field theoretic origin for the repulsive Dark Energy vacuum.
\end{abstract}

\maketitle

\section{Introduction}

The requirement that Effective Field Theories (EFTs) must couple consistently to quantum gravity places severe restrictions on low-energy phenomenology. Theories that fail these consistency conditions are relegated to the ``Swampland''. In recent years, Fuzzy Dark Matter (FDM)---composed of ultra-light scalar fields with masses $m_a \sim 10^{-21}$ eV---has emerged as a leading candidate to resolve the small-scale crises of Cold Dark Matter. 

In string phenomenology, the masses and couplings of scalar fields are rigidly determined by the topology and geometry of the internal compactified extra dimensions. In this work, we demonstrate that a Type IIB orientifold compactification on a K3-fibered Calabi-Yau threefold provides a mathematically motivated geometric origin for both the FDM scalar and the Dark Energy vacuum.

\begin{remark}[Epistemic Scope]
This paper establishes a phenomenological mapping between K3 topological invariants and the dark sector. We do not claim to have constructed a fully stabilized string vacuum. The moduli values required to achieve explicit mass scales are treated as free parameters, and exact tadpole flux vectors require future top-down derivation.
\end{remark}

\section{Topological Origins of the Dark Sector}

\subsection{Defining the Geometry}
In algebraic geometry, a K3 surface is uniquely characterized by three exact integer invariants: the second Betti number $b_2 = 22$, the Euler characteristic $\chi = 24$, and the Hirzebruch signature $\tau = b_2^+ - b_2^- = 3 - 19 = -16$. 

\subsection{Generate Dark Matter (The Axiverse): $b_2 = 22$}
Recent theoretical models, such as the 5D orbifold framework by Lee \& Tsai \cite{tsai2026}, propose that dark matter is a natural mass-resonance enforced by the geometry of extra dimensions. Our 6D framework provides a concrete string-theoretic realization of this mechanism.

Because a K3 surface has exactly $b_2 = 22$ independent 2-cycles, integrating the Ramond-Ramond 4-form ($C_4$) field over K3 mathematically yields exactly 22 distinct axion fields:
\begin{equation}
    a_i(x) = \int_{\Sigma_i} C_4, \quad i \in \{1, 2, \dots, 22\}.
\end{equation}
The masses of these axions are strictly generated by non-perturbative Euclidean D3-brane instantons wrapping the corresponding 4-cycles. The $S_{1,2}$ Picard-Fuchs sequence governs the specific topological rigidity of the primary cycle, defining the geometric resonance that sets the mass of the FDM axion to the required $10^{-21}$ eV, isolating it from the remaining spectator axions.

\subsection{Prove UV Consistency: Tadpole Cancellation ($\chi=24$)}
In F-theory and Type IIB orientifolds, the D3-brane charge must cancel against quantized fluxes to avoid fatal gauge anomalies:
\begin{equation}
    N_{D3} + \frac{1}{2\kappa_{10}^2 T_3} \int_{\mathcal{M}} H_3 \wedge F_3 = \frac{\chi}{24}.
\end{equation}
Because the Euler characteristic is strictly $\chi(\text{K3}) = 24$, the geometric contribution is exactly 1. This allows for an exact integer flux quantization yielding a stable, purely fluxed vacuum ($N_{D3} = 0$), mathematically avoiding the Swampland.

\subsection{Generate Dark Energy: $\chi=24$ and $\tau=-16$}
In the Large Volume Scenario (LVS), the non-zero Euler characteristic ($\chi = 24$) generates the perturbative $\alpha'^3$ quantum corrections to the K\"ahler potential. This stabilizes the volume of the extra dimensions while yielding a metastable de Sitter minimum (Dark Energy).

Furthermore, the Hirzebruch signature dictates the chiral asymmetry of the vacuum via the Atiyah-Singer Index Theorem ($n_+ - n_- = \tau = -16$). This severe chiral asymmetry prevents the perfect cancellation of vacuum fluctuations, sourcing the Trace Anomaly of the stress-energy tensor ($\langle \tensor{T}{^\mu_\mu} \rangle \neq 0$). This acts as an irreducible zero-point vacuum energy, establishing a geometric mechanism for a repulsive Dark Energy vacuum.

\begin{thebibliography}{99}
\bibitem{tsai2026} T. Lee and Y.-D. Tsai, \textit{Naturally resonant dark matter from extra dimensions}, Phys. Rev. D \textbf{114}, L011701 (2026).
\end{thebibliography}

\end{document}
